\chapter{Exercises and Reflection Prompts}
\label{ch:exercises}

\section{Short Answer Questions}

\begin{exercisebox}
\begin{exercise}
State the difference between framework inertia and volitional energy in one
paragraph. Why is it a mistake to collapse both into the single word
``motivation''?
\end{exercise}

\begin{exercise}
Give one example of a decision window that occurs in your own daily routine.
Explain why the window is short-lived.
\end{exercise}

\begin{exercise}
Describe a case in which the target state was too ambiguous to enter. Which
term in the barrier decomposition captured that ambiguity?
\end{exercise}
\end{exercisebox}

\section{Applied Modelling Tasks}

\begin{exercisebox}
\begin{exercise}
Model the transition from passive information consumption to active note-taking.
Name the current state, target state, dominant reward source, switching cost,
and one environmental reinforcement.
\end{exercise}

\begin{exercise}
Choose a repeated behaviour failure and design a one-page intervention using
the template from Chapter~\ref{ch:interventions}. Your answer must include a
decision window and a verification rule.
\end{exercise}
\end{exercisebox}

\section{Contrast and Boundary Questions}

\begin{exercisebox}
\begin{exercise}
Give an example of a case where the main problem is not high current inertia
but low clarity about long-run goals. Explain why the framework is still
useful, but not sufficient by itself.
\end{exercise}

\begin{exercise}
When might an external commitment lower barriers for one person but raise them
for another? Answer using the language of reward, environment, and cognitive
cost.
\end{exercise}
\end{exercisebox}

\section{Self-Audit Worksheet}

\begin{templatebox}
\textbf{Weekly framework inertia audit}
\begin{enumerate}[nosep]
  \item What beneficial transition did I fail to start this week?
  \item What was the current state that won instead?
  \item Which barrier term dominated the failure?
  \item Which decision window did I miss?
  \item What one preparation change will I install before the next attempt?
  \item What evidence next week would count as real improvement?
\end{enumerate}
\end{templatebox}

\section{Instructor-Oriented Prompting Uses}

This chapter can also support discussion-based teaching:
\begin{itemize}[nosep]
  \item ask students to compare two failed transitions with the same model;
  \item ask for a case where the best intervention is environmental rather than
        motivational;
  \item ask students to redesign a failure using a smaller entry action.
\end{itemize}

\section{Recommended Reading}

\begin{readingbox}
\textbf{Foundation}
\begin{itemize}[nosep]
  \item Revisit the notation appendix and the case library before attempting to
        extend the exercises mathematically.
\end{itemize}
\textbf{Modern Research}
\begin{itemize}[nosep]
  \item The literature map provides the best route for turning these exercises
        into research-backed problem sets later.
\end{itemize}
\textbf{Frontier}
\begin{itemize}[nosep]
  \item Frontier readings should be used for seminar prompts, not for baseline
        assessment.
\end{itemize}
\end{readingbox}

\begin{summarybox}
The exercise layer turns the manuscript from a reading artifact into a usable
training artifact. That matters because a textbook without exercises cannot
prove that its concepts transfer beyond the page.
\end{summarybox}
